\documentclass[a4paper,12pt]{article}
\usepackage{fontspec}
\usepackage{polyglossia}
\setmainlanguage{russian}
\setmainfont{DejaVu Serif}
\newfontfamily\cyrillicfont{DejaVu Serif}
\newfontfamily\cyrillicfonttt{DejaVu Sans Mono}
\usepackage{amsmath}
\usepackage{geometry}
\geometry{left=25mm,right=15mm,top=20mm,bottom=20mm}

\begin{document}

\begin{titlepage}
\begin{center}
МОСКОВСКИЙ АВИАЦИОННЫЙ ИНСТИТУТ\\
(НАЦИОНАЛЬНЫЙ ИССЛЕДОВАТЕЛЬСКИЙ УНИВЕРСИТЕТ)\\
Институт №8 «Компьютерные науки и прикладная математика»\\
Кафедра 806 «Вычислительная математика и программирование»
\vfill
{\Large Лабораторные работы №1--4}\\
по курсу «Численные методы»\\
Вариант №6
\vfill
\begin{flushright}
Выполнил: Ибрагимов Далгат Магомед алиевич\\
Группа: М8О-308Б-23\\
Преподаватель: Ревизников Д. Л.
\end{flushright}
\vfill
Москва, 2026
\end{center}
\end{titlepage}

\section*{Лабораторная работа №1}
Реализованы методы численной линейной алгебры:
\begin{itemize}
    \item LU-разложение с выбором главного элемента;
    \item метод прогонки;
    \item методы простых итераций и Зейделя;
    \item метод вращений Якоби;
    \item QR-алгоритм.
\end{itemize}
Исходные данные 6-го варианта находятся в файлах \texttt{lab1/lab1-* /input.txt}.
Результаты записаны в соответствующие файлы \texttt{output.txt}.

\section*{Лабораторная работа №2}
Для варианта 6 решены нелинейное уравнение и система нелинейных уравнений.
Использованы метод простых итераций и метод Ньютона. Код и результаты
находятся в каталогах \texttt{lab2/lab2-1} и \texttt{lab2/lab2-2}.

\section*{Лабораторная работа №3}
Реализованы интерполяция Лагранжа и Ньютона, кубический сплайн, метод
наименьших квадратов, численное дифференцирование и численное интегрирование
с оценкой погрешности методом Рунге--Ромберга. Данные 6-го варианта внесены
в файлы \texttt{input.txt}, результаты находятся в \texttt{output.txt}.

\section*{Лабораторная работа №4}
Реализованы методы решения задачи Коши и краевой задачи для ОДУ второго
порядка: метод Эйлера, метод Рунге--Кутты, метод Адамса, метод стрельбы и
конечно-разностный метод. Результаты для 6-го варианта сохранены в
\texttt{lab4/lab4-1/output.txt} и \texttt{lab4/lab4-2/output.txt}.

\end{document}
